\subsection{Федеральный закон № 187-ФЗ «О безопасности критической информационной инфраструктуры Российской Федерации»}

Федеральный закон № 187-ФЗ «О безопасности критической информационной инфраструктуры Российской Федерации» регулирует отношения, возникающие при обеспечении безопасности значимых объектов критической информационной инфраструктуры (КИИ). Закон распространяется на информационные системы, информационно-телекоммуникационные сети и автоматизированные системы управления, функционирование которых имеет ключевое значение для устойчивого функционирования государства, экономики и обеспечения безопасности граждан.

К субъектам КИИ относятся организации, осуществляющие деятельность в таких сферах, как здравоохранение, транспорт, связь, энергетика, банковская и финансовая сфера, топливно-энергетический комплекс, промышленность и иные отрасли. Для таких объектов устанавливаются специальные требования по обеспечению безопасности, включая проведение категорирования объектов КИИ, выявление актуальных угроз, реализацию мер защиты и организацию процессов реагирования на компьютерные инциденты.

В соответствии с положениями закона субъекты КИИ обязаны:
\begin{itemize}
    \item проводить категорирование объектов критической информационной инфраструктуры;
    \item обеспечивать реализацию организационных и технических мер по защите информации;
    \item осуществлять мониторинг состояния безопасности и выявление компьютерных инцидентов;
    \item обеспечивать взаимодействие с государственными системами обнаружения, предупреждения и ликвидации последствий компьютерных атак;
    \item принимать меры по предотвращению, обнаружению и ликвидации последствий компьютерных атак.
\end{itemize}

Поскольку Kubernetes-кластеры широко используются в современной корпоративной и государственной инфраструктуре, в том числе в организациях, относящихся к субъектам КИИ, данный закон имеет прямое отношение к теме выпускной квалификационной работы. Kubernetes-кластер может являться частью значимого объекта КИИ или обеспечивать функционирование критически важных сервисов.

Разрабатываемое средство обнаружения аномалий на основе анализа audit-логов Kubernetes-кластера может рассматриваться как элемент системы мониторинга безопасности значимого объекта КИИ. Оно позволяет осуществлять сбор и анализ событий, выявлять отклонения в поведении пользователей и сервисных учётных записей, обнаруживать признаки компрометации и несанкционированного доступа.

Таким образом, применение разработанного решения способствует выполнению требований Федерального закона № 187-ФЗ в части мониторинга безопасности, выявления и реагирования на инциденты, а также повышения общего уровня защищённости критической информационной инфраструктуры.